%%%%%%%%%%%%%%%%%%%%%%%%%%%%%%%%%%%%%%%%%%%%%%%%%%%%%%%
%% Unofficial class for ANR Project proposals
%%
%% v1.2.4 - 21/02/2023: updated content to match the 2023 version of the ANR template, thanks to Vincent Colotte.
%% v1.2.3 - 08/09/2021: adjusted footer and header, thanks to Adel Noureddine.
%% v1.2.2 - 21/11/2020: adjusted to allow lists in notes, thanks to Christophe Gravier.
%% v1.2.1 - 25/02/2020: modification to fit the changes in packages pgfgantt, thanks to Nicolas Marchand.
%% v1.2   - 02/03/2019: minor adjustments to match the 2019 version of the template.
%% v1.1   - 12/04/2018: minor adjustments in the class file.
%% v1     - 07/03/2018: first draft.
%%
%% Vincent Labatut 2018-23 
%% <vincent.labatut@univ-avignon.fr>
%% https://www.overleaf.com/latex/templates/unofficial-template-for-anr-proposals/yqgzsxkzrqkw
%% Creative Commons CC BY 4.0
%%
%% Please tell me if you find any error.
%% And if you use this model for your ANR project: 
%% I only take 8% of your accepted budget ;)
%%%%%%%%%%%%%%%%%%%%%%%%%%%%%%%%%%%%%%%%%%%%%%%%%%%%%%%
\documentclass{anr/proposal}





%%%%%%%%%%%%%%%%%%%%%%%%%%%%%%%%%%%%%%%%%%%%%%%%%%
% Search for "NOTE:" to get to the parameters worth adjusting to your needs
% There are some such NOTEs in the anr/proposal.cls file, too.
%%%%%%%%%%%%%%%%%%%%%%%%%%%%%%%%%%%%%%%%%%%%%%%%%%






%%%%%%%%%%%%%%%%%%%%%%%%%%%%%%%%%%%%%%%%%%%%%%%%%%
% project details
\acronym				        % acronym of the project
	{\textsc{MyProJecT}}
\committee				        % scientific evaluation committee (CES)
	{CES 48 -- Foundations of digital technology: computer science, automatic, signal processing}
\instrument{PRC}		        % funding instrument
\acyear{2023}			        % academic year of the call
\title                          % full title of the project
    {\myacronym{}: Proposal's title}
\subtitle                           % subtitle of the project (optional)
    {Whatever subtitle providing additional information}
\investigator{Stéphanie MARTIN} % principal investigator
\duration{36 months}            % duration of the project
\funding{917 637\euro}          % requested funding
% NOTE: all the above fields must be filled (except the subtitle which is optional)
% NOTE: use command \myxxxxx{} to mention information xxxx in the text, e.g. \myacronym{}


%%%%%%%%%%%%%%%%%%%%%%%%%%%%%%%%%%%%%%%%%%%%%%%%%%
% bibliography
\addbibresource{anr_biblio.bib}
% NOTE: look for "bibliographic settings" in the class if you want to 
% change the formatting of the bibliography (e.g. switch to author-year instead of numeric)


%%%%%%%%%%%%%%%%%%%%%%%%%%%%%%%%%%%%%%%%%%%%%%%%%%
% partners
\newcommand{\colUE}[1]{\textcolor{partnercol1}{#1}}
\newcommand{\colEGS}[1]{\textcolor{partnercol2}{#1}}
\newcommand{\colINR}[1]{\textcolor{partnercol3}{#1}}
\newcommand{\colUT}[1]{\textcolor{partnercol4}{#1}}
\newcommand{\colTUT}[1]{\textcolor{partnercol5}{#1}}
% NOTE: colors go from partnercol1 to partnercol8
% if you need more colors, add them in the class file, look for "color settings"


%%%%%%%%%%%%%%%%%%%%%%%%%%%%%%%%%%%%%%%%%%%%%%%%%%
% work packages
% NOTE: WP are defined using the \wpdef{label}{Title} macro. 
% NOTE: Cross-referencing works as usual, with the \ref{label} macro.
% NOTE: WPs are numbered from 0. If you want to start from 1, look for "work packages and tasks" 
% in the class file (anr.cls), and read the note there.


%%%%%%%%%%%%%%%%%%%%%%%%%%%%%%%%%%%%%%%%%%%%%%%%%%
% tasks
% NOTE: tasks are defined using the \tdef{label}{Title} macro. 
% NOTE: Cross-referencing works as usual, with the \ref{label} macro.


%%%%%%%%%%%%%%%%%%%%%%%%%%%%%%%%%%%%%%%%%%%%%%%%%%
% writer comments
\usepackage%[disable]	% uncomment to hide all the notes
	{todonotes} 
\newcommand{\noteANR}[1]{\todo[color=gray!40, author=\textbf{ANR}, inline, caption={}]{#1}}
\newcommand{\noteSM}[1]{\todo[color=red!40, author=\textbf{Stéphanie}, inline, caption={}]{#1}}
\newcommand{\noteJZ}[1]{\todo[color=green!40, author=\textbf{Jing}, inline, caption={}]{#1}}
\newcommand{\noteRM}[1]{\todo[color=blue!40, author=\textbf{Richard}, inline, caption={}]{#1}}
% NOTE: you can setup other writers by pasting and editing the above lines


\begin{document}
\maketitle




%%%%%%%%%%%%%%%%%%%%%%%%%%%%%%%%%%%%%%%%%%%%%%%%%%
%%%%%%%%%%%%%%%%%%%%%%%%%%%%%%%%%%%%%%%%%%%%%%%%%%
%%%%%%%%%%%%%%%%%%%%%%%%%%%%%%%%%%%%%%%%%%%%%%%%%%
%%%%%%%%%%%%%%%%%%%%%%%%%%%%%%%%%%%%%%%%%%%%%%%%%%




%%%%%%%%%%%%%%%%%%%%%%%%%%%%%%%%%%%%%%%%%%%%%%%%%%
% consortium
\begin{center}
	\textbf{Consortium:} 
	\colUE{LITA (UE)}, 
	\colEGS{LMQ (EGS)},
	\colINR{LMD (INR)},
	\colUT{CGCS (UT)},
	\colTUT{FI (TUT)}
\end{center}




%%%%%%%%%%%%%%%%%%%%%%%%%%%%%%%%%%%%%%%%%%%%%%%%%%
% general comments
\noteANR{Please use an easily readable document layout (A4 pages, Calibri 11 or equivalent, single spaced, 2cm margins, numbered pages ; for figure and table, minimum Calibri 9 or equivalent).}

\noteANR{The project description (1) cannot exceed a 20-page limit (including summary table of persons involved in the project, Gantt chart, overview of the implication of scientific leaders in on-going projects, overview of the requested funds AND their scientific justification, and bibliography) and (2) must be submitted in a PDF format. The use of hypertext links leading to the bibliography of the project, figures, CVs, among other examples, amounts to not respecting the limit of 20 pages applicable to all coordinators.}
\noteANR{CVs of the scientific coordinator and any partners’ scientific leaders must be completed online, on \href{https://iris.anr.fr/fr/login}{IRIS}, before the closing date of stage 2: each scientific leader must login to the Website \href{https://iris.anr.fr/fr/login}{IRIS}, and click on his/her name, on the top-right of the homepage to fulfill his/her resume}

\noteANR{As the evaluation could be carried out by non-French speaking researchers, we recommend that you write both your scientific document and the CVs in English.}

\noteANR{Proposals must fullfil the three main evaluation criteria: “Quality and scientific aims”, “Organisation and implementation of the project”, “Impact and benefits of the project”. Applicants are advised to consult the document \href{https://anr.fr/fileadmin/aap/2023/aapg-20203-Guide_V2.0_20220921.pdf}{AAPG 2023 Guide} for further information about the different sub-criteria related to the chosen funding instrument.}

\noteANR{The review of the projects is based on the information as completed and submitted on \href{https://iris.anr.fr/fr/login}{IRIS} at the closing date and time (within the limit of 20 pages concerning the scientific document). No information other than that completed and submitted at the closing date and time on \href{https://iris.anr.fr/fr/login}{IRIS} will be requested from the coordinators or searched on the Internet.}

\noteANR{\textbf{Evaluation criteria for projects submitted to the Generic Call for Proposals 2019 :}}

\noteANR{\textbf{1) Quality and scientific aims} (Phase 2): \\
-- Clarity of research objectives and hypotheses. \\
-- Scientific ambition and position in relation to the state of the art. \\
-- Appropriateness of the methodology, management of scientific risks.}

\noteANR{\textbf{2) Organisation and implementation of the project} (Phase 2): \\
-- Competences, expertise and involvement of the scientific coordinator and the partners. \\
-- For \textbf{PRC/PRCE/PRCI} proposals: Quality and complementarity of the consortium,
quality of collaboration, or For \textbf{JCJC} proposals: Contribution to the coordinator’s level of responsibility and team development. \\
-- Appropriateness of implemented and requested resources to the project’s objectives.}

\noteANR{\textbf{3) Impact and benefits of the project:} \\ 
-- Scientific impact and potential economic, social or cultural impact
-- For \textbf{PRC/JCJC} proposals: Research outcome dissemination and application strategy,
including promoting scientific culture, or For \textbf{PRCE} proposals: Action of technology and innovation transfer with regard to the social and economic world, or For \textbf{PRCI} proposals: Balance and complementarity of scientific contributions by respective partners in each country and added value/benefit for France from this European or international cooperation.}




%%%%%%%%%%%%%%%%%%%%%%%%%%%%%%%%%%%%%%%%%%%%%%%%%%
% project summary
\input{anr_0_summary.tex}




%%%%%%%%%%%%%%%%%%%%%%%%%%%%%%%%%%%%%%%%%%%%%%%%%%
% Context and objectives
\input{anr_1_context.tex}




%%%%%%%%%%%%%%%%%%%%%%%%%%%%%%%%%%%%%%%%%%%%%%%%%%
\input{anr_2_organization.tex}




%%%%%%%%%%%%%%%%%%%%%%%%%%%%%%%%%%%%%%%%%%%%%%%%%%
\input{anr_3_impact.tex}




%%%%%%%%%%%%%%%%%%%%%%%%%%%%%%%%%%%%%%%%%%%%%%%%%%
\section{References related to the project}

\noteANR{This paragraph refers to the evaluation criteria « “Quality and scientific aim”}

\noteANR{List the bibliographical references used for the proposal.}

\noteANR{Please, fill in “usable” references, i.e. including the first co-authors, complete title, title of the journal, etc. If available, please complete these references (but not replace) by indicating the « open access » link to improve accessibility for the reviewers.}

\noteANR{Preprints are allowed, especially those referencing preliminary data.}

\noteANR{Impact factors are prohibited.}

\noteANR{\textbf{The bibliography must be included in the 20-page limit.}}

% Bibliography
{	\setstretch{0.8}							% decreases the space between lines
	\phantomsection
	\addcontentsline{toc}{section}{References}	% add bibliography to toc
	\printbibliography							% add the references
}

\end{document}




%%%%%%%%%%%%%%%%%%%%%%%%%%%%%%%%%%%%%%%%%%%%%%%%%%
%%%%%%%%%%%%%%%%%%%%%%%%%%%%%%%%%%%%%%%%%%%%%%%%%%
%%%%%%%%%%%%%%%%%%%%%%%%%%%%%%%%%%%%%%%%%%%%%%%%%%
% TODO
% garder les sections* dans les liens du PDF
